\section{数据整理与合并}

大多数真正的工作从这里开始。你通常面对两个或更多数据集，主键不明确，朴素的 merge 可能丢数据或产生重复行。

\subsection{合并前的 Profiling}

要求 Codex 在执行合并\textbf{之前}先做以下检查：

\begin{enumerate}[nosep]
    \item \textbf{检查候选键的唯一性}（uniqueness of candidate keys）
    \item \textbf{度量空值率和格式差异}（null rates and formatting differences）
    \item \textbf{规范化明显的格式问题}：大小写、空格、地址格式等
    \item \textbf{执行试验性 join 并报告匹配率}（trial joins and match rates）
    \item \textbf{推荐最安全的合并策略}，然后才写入最终合并文件
\end{enumerate}

\begin{importantbox}{``先 Profile，再 Merge''原则}
如果需要构造最优的 join key（如规范化地址、从多列拼接的地块标识符、或基于位置的 join），让 Codex 先解释 trade-off 和边界情况，\textbf{你确认后再执行合并}。这是防止数据泄漏和合并错误的关键步骤。
\end{importantbox}

\subsection{本章小结}
数据合并是分析中最容易出错的环节。核心原则是``先审查、再合并''：让 Codex 报告候选键、空值率、匹配率和策略建议，确认无误后才执行。

